% !TEX root = main.tex
% Math - some things from physics package need to be redefined with \sym*f for compatibility with unicode-math and accessible pdfs.
% For commands rewritten by some packages, useful to have \AtBeginDocument to ensure implementation actually occurs
% \AtBeginDocument{\renewcommand{\mathrm}[1]{\symup{#1}}}
\AtBeginDocument{\renewcommand{\mathbf}[1]{\symbf{#1}}}
\newcommand{\diffd}{\symup{d}}
\AtBeginDocument{\renewcommand{\div}{ \symbf{\nabla} \cdot }}
\newcommand{\grad}{ \symbf{\nabla} }
\newcommand{\gradient}{ \symbf{\nabla} }
\newcommand{\curl}{ \symbf{\nabla} \times } 
\newcommand{\vdot}{ \cdot }

\DeclareDocumentCommand\crossproduct{}{\symbf\times} % Vector cross product symbol
\DeclareDocumentCommand\cross{}{\crossproduct} % Shorthand for \crossproduct
\DeclareDocumentCommand\cp{}{\crossproduct} % Shorthand for \crossproduct

\DeclareMathOperator{\sech}{sech}
% Unit vector
\DeclareDocumentCommand\vectorunit{ s m }{\IfBooleanTF{#1}{\boldsymbol{\hat{#2}}}{\symbf{\hat{#2}}}} % Unit vector [star for Greek and italic Roman]
\DeclareDocumentCommand\vu{}{\vectorunit} % Shorthand for \vectorunit


\DeclareDocumentCommand\vectorbold{ s m }{\IfBooleanTF{#1}{\boldsymbol{#2}}{\symbf{#2}}} % Vector bold [star for Greek and italic Roman]
\DeclareDocumentCommand\vb{}{\vectorbold} % Shorthand for \vectorbold

% \AtBeginDocument{\newcommand{\div}{ \symbf{\nabla} \cdot } }
% \AtBeginDocument{\newcommand{\grad}{ \symbf{\nabla} } }
% \AtBeginDocument{\newcommand{\curl}{ \symbf{\nabla} \times } }
% \AtBeginDocument{\newcommand{\vdot}{ \cdot } }

% \AtBeginDocument{\newcommand{\mathbf}[1]{\symbf{#1}}}

%derivative command from physics package 
\DeclareDocumentCommand\derivative{ s o m g d() }
{ % Total derivative
	% s: star for \flatfrac flat derivative
	% o: optional n for nth derivative
	% m: mandatory (x in df/dx)
	% g: optional (f in df/dx)
	% d: long-form d/dx(...)
	\IfBooleanTF{#1}
	{\let\fractype\flatfrac}
	{\let\fractype\frac}
	\IfNoValueTF{#4}
	{
		\IfNoValueTF{#5}
		{\fractype{\diffd \IfNoValueTF{#2}{}{^{#2}}}{\diffd #3\IfNoValueTF{#2}{}{^{#2}}}}
		{\fractype{\diffd \IfNoValueTF{#2}{}{^{#2}}}{\diffd #3\IfNoValueTF{#2}{}{^{#2}}} \argopen(#5\argclose)}
	}
	{\fractype{\diffd \IfNoValueTF{#2}{}{^{#2}} #3}{\diffd #4\IfNoValueTF{#2}{}{^{#2}}}}
}
\DeclareDocumentCommand\dv{}{\derivative} % Shorthand for \derivative

% Partial Derivative command from phyiscs
\DeclareDocumentCommand\partialderivative{ s o m g g d() }
{ % Partial derivative
	% s: star for \flatfrac flat derivative
	% o: optional n for nth derivative
	% m: mandatory (x in df/dx)
	% g: optional (f in df/dx)
	% g: optional (y in d^2f/dxdy)
	% d: long-form d/dx(...)
	\IfBooleanTF{#1}
	{\let\fractype\flatfrac}
	{\let\fractype\frac}
	\IfNoValueTF{#4}
	{
		\IfNoValueTF{#6}
		{\fractype{\partial \IfNoValueTF{#2}{}{^{#2}}}{\partial #3\IfNoValueTF{#2}{}{^{#2}}}}
		{\fractype{\partial \IfNoValueTF{#2}{}{^{#2}}}{\partial #3\IfNoValueTF{#2}{}{^{#2}}} \argopen(#6\argclose)}
	}
	{
		\IfNoValueTF{#5}
		{\fractype{\partial \IfNoValueTF{#2}{}{^{#2}} #3}{\partial #4\IfNoValueTF{#2}{}{^{#2}}}}
		{\fractype{\partial^2 #3}{\partial #4 \partial #5}}
	}
}
\DeclareDocumentCommand\pderivative{}{\partialderivative} % Shorthand for \partialderivative
\DeclareDocumentCommand\pdv{}{\partialderivative} % Shorthand for \partialderivative

\DeclareDocumentCommand\poissonbracket{ l m m }{\braces#1{\lbrace}{\rbrace}{#2,#3}} % Poisson bracket [same as anti-commutator]
\DeclareDocumentCommand\pb{}{\poissonbracket} % Shorthand for \poissonbracket
